\section*{BIBLIOGRAPHY}
\phantomsection
\addcontentsline{toc}{section}{Bibliography}

\begin{enumerate}
\def\labelenumi{\arabic{enumi}.}
\item
  O. NEUGEBAUER, Vorlesungen über die Geschichte der antiken Mathematik, Vol. I, Vorgriechische Mathematik, Berlin (Springer), 1934.
\item
  The Works of Aristotle, translated under the editorship of J. A. Smith and W. D. Ross (12 volumes, Oxford, 1908-1952).
\end{enumerate}

2 (bis). T. L. HEATH, Mathematics in Aristotle, Oxford (Clarendon Press), 1949.

\begin{enumerate}
\def\labelenumi{\arabic{enumi}.}
\setcounter{enumi}{2}
\item
  T. L. HEATH, The Thirteen Books of Euclid's Elements\ldots, 3 volumes, Cambridge, 1908.
\item
  Archimedis Opera Omnia, 3 volumes, edited by J. L. Heiberg, 2nd edition, 1913-1915.
\end{enumerate}

4 (bis). T. L. HEATH, The Method of Archimedes, Cambridge, 1912.

\begin{enumerate}
\def\labelenumi{\arabic{enumi}.}
\setcounter{enumi}{4}
\item
  J. M. BOCHENSKI, Ancient Formal Logic, Studies in Logic, Amsterdam (North Holland Publ. Co.), 1951.
\item
  P. BÖHNER, Medieval Logic, an Outline of its Development from 1250 to ca. 1400, Chicago, 1952.
\item
  D. FRANCISCI MAUROLYCI, Abbatis Messamensis, Mathematici Celeberrimi, Arithmeticorum Libri Duo, Venice, 1575.
\item
  GALILEO GALILEI, Opere, Ristampa della Edizione Nazionale, 20 volumes, Florence (Barbara), 1929-39.
\end{enumerate}

(*) I. NOVAK-GAL, ``A construction for models of consistent systems'', Fund. Math., 37 (1950), pp.~87-110.

\begin{enumerate}
\def\labelenumi{\arabic{enumi}.}
\setcounter{enumi}{8}
\item
  A. GIRARD, Invention nouvelle en l'Algèbre, 1629 (new edition, Bierens de Haan, 1884).
\item
  R. DESCARTES, Œuvres, edited by C. Adam and P. Tannery, 11 volumes, Paris (L. Cerf), 1897-1909.
\item
  B. PASCAL, Œuvres, edited by Brunschvicg, 14 volumes, Paris (Hachette), 1904-1914.
\item
  G. W. LEIBNIZ : (a) Mathematische Schriften, edited by C. I. Gerhardt, 7 volumes, Berlin-Halle (Asher-Schmidt), 1849-1863; (b) Philosophische Schriften, edited by C. I. Gerhardt, 7 volumes, Berlin, 1840-1890; (c) Opuscules et fragments inédits, edited by L. Couturat, Paris (Alcan), 1903.
\end{enumerate}

12 (bis). L. COUTURAT, La logique de Leibniz d'après des documents inédits, Paris (Alcan), 1901.

\begin{enumerate}
\def\labelenumi{\arabic{enumi}.}
\setcounter{enumi}{12}
\item
  D'ALEMBERT, Encyclopédie, Paris, 1751-1765, articles ``Négatif'', ``Imaginaire'', ``Définition''.
\item
  C. F. GAUSS, Werke, 12 volumes, Göttingen, 1870-1927.
\item
  N. LOBATSCHEVSKY, Pangeometrie, Ostwald's Klassiker, no. 130, Leipzig (Engelmann), 1902.
\item
  G. BOOLE : (a) The Mathematical Analysis of Logic, Cambridge-London, 1847 (= Collected Logical Works, edited by P. Jourdain, Chicago-London, 1916, vol.~I); (b) An Investigation of the Laws of Thought, Cambridge-London, 1854 (= Collected Logical Works, vol.~II).
\item
  H. GRASSMANN, Gesammelte Werke, 6 volumes, Leipzig (Teubner), 1894.
\item
  B. BOLZANO, Paradoxien des Unendlichen, Leipzig, 1851.
\item
  B. RIEMANN, Gesammelte mathematische Werke, 2nd edition, Leipzig (Teubner), 1892.
\item
  H. HANKEL, Theorie der complexen Zahlensysteme, Leipzig (Voss), 1867.
\item
  A. DE MORGAN : (a) ``On the syllogism (III)'', Trans. Camb. Phil. Soc., 10 (1858), pp.~173-230; (b) ``On the syllogism (IV) and on the logic of relations'', Trans. Camb. Phil. Soc., 10 (1860) pp.~331-358.
\item
  C. S. PEIRCE, (a) ``Upon the logic of mathematics'', Proc. Amer. Acad. Arts and Sci., 7 (1865-1868), pp.~402-412; (b) ``On the algebra of logic'', Amer. Journ. of Math., 3 (1880), pp.~49-57; (c) ``On the algebra of logic'', Amer. Journ. of Math., 7 (1884), pp.~190-202.
\item
  E. SCHRÖDER, Vorlesungen über die Algebra der Logik, 3 volumes, Leipzig (Teubner), 1890.
\item
  G. FREGE, (a) Begriffsschrift eine der arithmetischen nachgebildete Formelsprache des reinen Denkens, Halle, 1879; (b) Die Grundlagen der Arithmetik, 2nd edition with an English translation by J. L. Austin, New York, 1950; (c) Grundgesetze der Arithmetik, begriffsschriftlich abgeleitet, 2 volumes, Jena, 1893-1903.
\item
  G. CANTOR, Gesammelte Abhandlungen, Berlin (Springer), 1932.
\end{enumerate}

25 (bis). G. CANTOR, R. DEDEKIND, Briefwechsel, edited by J. Cavaillès and E. Noether, Act. Sci. et Ind., no. 518, Paris (Hermann), 1937.

\begin{enumerate}
\def\labelenumi{\arabic{enumi}.}
\setcounter{enumi}{25}
\item
  R. DEDEKIND, Gesammelte mathematische Werke, 3 volumes, Braunschweig (Vieweg), 1932.
\item
  C. HERMITE, T. STIELTJES, Correspondance, 2 volumes, Paris (Gauthier-Villars), 1905.
\item
  M. PASCH und M. DEHN, Vorlesungen über neuere Geometrie, 2nd edition, Berlin, (Springer), 1926.
\item
  G. PEANO, (a) Arithmeticae principia, novo modo exposita, Turin, 1889; (b) I principii di Geometria, logicamente espositi, Turin, 1889; (c) Formulaire de Mathematiques, 5 volumes, Turin, 1895-1905.
\item
  D. HILBERT, Grundlagen der Geometrie, 7th edition, Leipzig-Berlin (Teubner), 1930.
\item
  D. HILBERT, Gesammelte Abhandlungen, vol.~III, Berlin (Springer), 1935.
\item
  D. HILBERT, W. ACKERMANN, Grundzüge der theoretischen Logik, 3rd edition, Berlin (Springer), 1949.
\item
  H. POINCARÉ: (a) Science et hypothèse, Paris (Flammarion), 1906; (b) La valeur de la science, Paris (Flammarion), 1905; (c) Science et méthode, Paris (Flammarion), 1920.
\item
  J. RICHARD, ``Les principes des mathématiques et le problème des ensembles'', Rev.~Gen.~des Sci. Pures et Appl., 16 (1905), pp.~541-543.
\item
  R. BAIRE, E. BOREL, J. HADAMARD, H. LEBESGUE, ``Cinq lettres sur la théorie des ensembles'', Bull. Soc. Math. de France, 33 (1905), pp.~261-273.
\item
  E. ZERMELO, (a) ``Beweis dass jede Menge wohlgeordnet werden kann'', Math. Ann., 59 (1904), pp.~514-516; (b) ``Neuer Beweis für die Möglichkeit einer Wohlordnung'', Math. Ann., 65 (1908), pp.~107-128; (c) ``Untersuchungen über die Grundlagen der Mengenlehre'', Math. Ann., 65 (1908), pp.~261-281.
\item
  B. RUSSELL and A. N. WHITEHEAD, Principia Mathematica, 3 volumes, Cambridge, 1910-1913.
\item
  L. E. J. BROUWER, (a) ``Intuitionism and formalism'', Bull. Amer. Math. Soc., 20 (1913), pp.~81-96; (b) ``Zur Begründung des intuitionistischen Mathematik'', Math. Ann., 93 (1925), pp.~244-257; 95 (1926), pp.~453-473; 96 (1926), pp.~451-458.
\item
  T. SKOLEM, ``Einige Bemerkungen zur axiomatischen Begründung der Mengenlehre'', Wiss. Vorträge, 5 Kongress Skand. Math., Helsingfors, 1922.
\item
  K. KURATOWSKI, ``Une méthode d'élimination des nombres transfinis des raisonnements mathématiques'', Fund. Math., 5 (1922), pp.~76-108.
\item
  A. FRAENKEL: (a) ``Zu den Grundlagen der Cantor-Zermeloschen Mengenlehre'', Math. Ann., 86 (1922), pp.~230-237; (b) Zehn Vorlesungen über die Grundlegung der Mengenlehre, Wiss. und Hypothese vol.~31, Leipzig-Berlin, 1927; (c) Einleitung in die Mengenlehre, 3rd edition, Berlin (Springer), 1928.
\item
  J. HERBRAND, ``Recherches sur la théorie de la démonstration'', Trav. Soc. Sci. et Lett. Varsovie, cl. II (1930), pp.~33-160.
\item
  J. VON NEUMANN: (a) ``Eine Axiomatisierung der Mengenlehre'', Crelle, 154 (1925), pp.~219-240; (b) ``Die Axiomatisierung der Mengenlehre'', Math. Zeitschr., 27 (1928), pp.~669-752; (c) ``Zur Hilbertschen Beweistheorie'', Math. Zeitschr., 26 (1927), pp.~1-46.
\item
  K. Gödel: (a) ``Über formal unentscheidbare Sätze der Principia Mathematica und verwandter Systeme'', Monatsh. für Math. u. Phys., 38 (1931), pp.~173-198; (b) The consistency of the axiom of choice and of the generalized continuum hypothesis, Ann. of Math. Studies, no. 3, Princeton, 1940.
\item
  M. ZORN, ``A remark on method in transfinite algebra'', Bull. Amer. Math. Soc., 41 (1935), pp.~667-670.
\item
  A. HEYTING, Mathematische Grundlagenforschung. Intuitionismus. Beweistheorie, Ergebnisse der Math., vol.~3, Berlin (Springer), 1934.
\item
  G. GENTZEN, Die gegenwärtige Lage in der mathematischen Grundlagenforschung. Neue Fassung des Widerspruchsfreiheitsbeweises für die reine Zahlentheorie, Forschungen der Logik\ldots, Heft 4, Leipzig (Hirzel), 1938.
\item
  S. KLEENE, Introduction to Metamathematics, New York, 1952.
\item
  P. J. COHEN, ``The independence of the continuum hypothesis'', Proc. Nat. Acad. Sci., 50 (1963), pp.~1143-1148 and 51 (1964), pp.~105-110.
\end{enumerate}
